\subsection{Exact anchor diversity--repetition pilot}
A subsequent exact-schedule pilot separated unique paired anchors $N_p$ from presentations per anchor $R_p$. Every selected anchor appeared exactly $R_p$ times in a deterministic shuffled schedule. The pilot crossed $N_u\in\{750,1500\}$, $N_p\in\{50,125,225\}$, $R_p\in\{32,64,128\}$, two methods, and five new seeds, for 180 neural-network fits.

Increasing repetition improved the universal biological score in all 12 fixed-$(N_u,N_p,\text{method})$ slices. Increasing unique-anchor diversity also improved it in all 12 fixed-$(N_u,R_p,\text{method})$ slices.

\begin{table}[t]
\centering
\caption{Universal biological score at $N_p=225$ and $R_p=128$.}
\small
\begin{tabular}{lcc}
\toprule
Method & $N_u=750$ & $N_u=1500$\\
\midrule
Hybrid curriculum & 0.503 & 0.524\\
Pair consistency & 0.495 & 0.505\\
\bottomrule
\end{tabular}
\end{table}

Task AUC remained approximately $0.96$--$0.98$ while biological recovery, factor separation, and paired-invariance error changed substantially. Strong task prediction therefore did not certify a disentangled biological representation.

\subsection{Matched paired-supervision budgets}
To separate unique intervention diversity from repeated exposure, two exact-schedule experiments fixed total paired presentations within each budget while changing the allocation between $N_p$ and $R_p$. The first compared
\begin{equation}
50\times128=100\times64=200\times32=6400,
\end{equation}
with 100 pair-loss-active steps. The second doubled repetition at every anchor count,
\begin{equation}
50\times256=100\times128=200\times64=12800,
\end{equation}
with 200 pair-loss-active steps. Both experiments used the same 10 matched seeds, $N_u\in\{750,1500\}$, hybrid curriculum and pair consistency, 320 optimizer steps, and 40,960 observational presentations. Each budget therefore comprised 120 neural-network fits.

\begin{table}[t]
\centering
\caption{Mean recovery across four method-by-$N_u$ strata under matched paired-supervision allocations.}
\small
\begin{tabular}{rccc}
\toprule
Budget & $N_p\times R_p$ & Biological score & Factor separation\\
\midrule
6400 & $50\times128$ & 0.4081 & 0.7487\\
6400 & $100\times64$ & 0.4243 & 0.7626\\
6400 & $200\times32$ & 0.4259 & 0.7661\\
12800 & $50\times256$ & 0.4489 & 0.7874\\
12800 & $100\times128$ & 0.4595 & 0.7922\\
12800 & $200\times64$ & 0.4619 & 0.7987\\
\bottomrule
\end{tabular}
\end{table}

For global allocation contrasts, the four method-by-$N_u$ cells were first averaged within each seed, leaving 10 independent matched seed blocks. At 6400 presentations, 200 anchors exceeded 50 anchors by 0.0179 (95\% seed-bootstrap interval $[0.0106,0.0261]$), with a positive difference in all 10 seed blocks and an exact two-sided sign-flip $p=0.00195$. The 100-versus-50 contrast was also positive on average, but its interval included zero ($0.0163$, interval $[-0.0017,0.0328]$, $p=0.113$). At 12800 presentations, the 200-versus-50 contrast remained positive ($0.0131$, interval $[0.0015,0.0251]$), although the exact test was less conclusive ($p=0.0762$). The 100-versus-50 contrast was $0.0106$ ($p=0.195$).

Differences between 200 and 100 anchors were near zero at both budgets: 0.0016 at 6400 presentations and 0.0024 at 12800 presentations, with exact $p$-values of 0.822 and 0.607. The allocation result therefore indicates that concentrating all supervision on only 50 unique anchors is inefficient in this regime, while the incremental benefit beyond approximately 100 anchors is small and uncertain.

Doubling total paired exposure produced the larger and more reproducible effect. Averaged within each seed across all 12 matched cells, 12800 presentations improved the universal biological score over 6400 by 0.0374 (95\% seed-bootstrap interval $[0.0303,0.0437]$). All 10 seed blocks improved and the exact two-sided sign-flip test gave $p=0.00195$. The improvement was positive for each allocation separately: 0.0408 at 50 anchors, 0.0352 at 100 anchors, and 0.0360 at 200 anchors.

These controlled results identify total paired exposure as the dominant resource over this range, with a smaller diversity-allocation effect and diminishing returns between 100 and 200 unique anchors. No allocation mean at either budget crossed the predefined recovery threshold of 0.50, so the experiments characterize resource efficiency and response to supervision rather than a confirmed recovery boundary.
